\documentclass[12pt,a4paper]{article}
\usepackage[utf8]{inputenc}
\usepackage[french]{babel}
\usepackage{amsmath}
\usepackage{geometry}
\usepackage{listings}
\usepackage{xcolor}
\usepackage{fancyhdr}
\usepackage{hyperref}

\geometry{margin=2.5cm}

\lstset{
  basicstyle=\ttfamily\small,
  keywordstyle=\color{red},
  frame=single,
  breaklines=true,
  showstringspaces=false
}

\pagestyle{fancy}
\fancyhf{}
\rhead{Recherche Opérationnelle 2025--2026}
\lhead{Université  -- Mauritanie}
\cfoot{\thepage}

\begin{document}

\begin{titlepage}
  \centering
  \vspace*{2cm}
  {\Large \textbf{Université mauritanie}\\
  Faculté des Sciences technique\\[0.5cm]}
  \rule{\linewidth}{0.5mm}\\[1cm]
  {\Huge \textbf{Rapport de Projet}}\\[0.5cm]
  {\LARGE Recherche Opérationnelle}\\[0.3cm]
  {\large Licence -- Année 2025--2026}\\[1cm]
  \rule{\linewidth}{0.5mm}\\[1cm]
  {\large \textbf{Sujets traités :}}\\[0.5cm]
  {\large
  \textbf{1.3} Problème du Sac à dos (Knapsack)\\[0.3cm]
  \textbf{1.1} Optimisation du Transport\\[0.3cm]
  \textbf{6.1} Transport de Poissons (Nouadhibou -- Nouakchott)\\
  }
  \vfill
  {\large
  \textbf{Responsable :} Dr. EL BENANY Med Mahmoud\\[0.3cm]
  \textbf{Étudiant :} [abdoulay hamet sy\\ matricule:c26058]\\[0.3cm]
  \textbf{Date :} 18 Mai 2026
  }
\end{titlepage}

\tableofcontents
\newpage

%==============================================
\section{Introduction}
%==============================================

La Recherche Opérationnelle est une discipline qui utilise les
mathématiques et les méthodes d'optimisation pour résoudre des
problèmes complexes.

Dans ce projet, nous avons utilisé Python pour traiter trois
problèmes principaux :
\begin{itemize}
  \item le problème du Sac à dos avec la programmation dynamique ;
  \item l'optimisation du transport de marchandises ;
  \item le transport de poissons frais en Mauritanie.
\end{itemize}

L'objectif général est de maximiser la valeur, minimiser les coûts
et optimiser les ressources disponibles.

Tous les codes sont disponibles sur GitHub :
\url{https://github.com/papasy1336-hub/projet-RO-2026}

%==============================================
\section{Problème du Sac à dos (Knapsack)}
%==============================================

\subsection{Description}
Le fichier \texttt{knapsack.py} traite le problème du sac à dos.
L'objectif est de sélectionner des objets pour maximiser la valeur
totale tout en respectant une contrainte de capacité maximale.

\subsection{Extrait du code}

\begin{lstlisting}
x = LpVariable.dicts("x", objets, cat="Binary")
\end{lstlisting}

\subsection{Modèle mathématique}

Le modèle utilisé est la programmation linéaire en nombres entiers.
Les variables de décision sont :

\[
x_i = \begin{cases} 1 & \text{si l'objet } i \text{ est sélectionné} \\ 0 & \text{sinon} \end{cases}
\]

\subsection{Fonction objectif}

La fonction objectif est :
\[
\max Z = 3x_1 + 4x_2 + 5x_3
\]

Dans le programme :
\begin{lstlisting}
modele += lpSum(valeurs[i] * x[i] for i in objets)
\end{lstlisting}

\subsection{Contraintes}

La contrainte de capacité est :
\[
2x_1 + 3x_2 + 4x_3 \leq 5
\]

Dans le programme :
\begin{lstlisting}
modele += lpSum(poids[i] * x[i] for i in objets) <= capacite
\end{lstlisting}

\subsection{Résultat obtenu}

\begin{lstlisting}
Statut : Optimal
Objet 1 (poids=2, valeur=3) -> PRIS
Objet 2 (poids=3, valeur=4) -> PRIS
Objet 3 (poids=4, valeur=5) -> non pris

Valeur totale = 7
\end{lstlisting}

La solution optimale permet de maximiser la valeur totale à
\textbf{7} avec un poids total de 5 (= capacité maximale).

%==============================================
\section{Optimisation du Transport}
%==============================================

\subsection{Description}
Le fichier \texttt{transport.py} traite le problème classique
de transport. L'objectif est de transporter des marchandises
entre plusieurs sources et plusieurs destinations avec un coût
minimal.

\subsection{Variables de décision}

Les variables $x_{ij}$ représentent la quantité transportée
de la source $i$ vers la destination $j$.

\begin{lstlisting}
x11 = LpVariable('x11', lowBound=0)
x12 = LpVariable('x12', lowBound=0)
x13 = LpVariable('x13', lowBound=0)

x21 = LpVariable('x21', lowBound=0)
x22 = LpVariable('x22', lowBound=0)
x23 = LpVariable('x23', lowBound=0)
\end{lstlisting}

\subsection{Fonction objectif}

La fonction objectif est :
\[
\min Z = 2x_{11} + 3x_{12} + x_{13} + 5x_{21} + 4x_{22} + 8x_{23}
\]

Dans le programme :
\begin{lstlisting}
modele += lpSum(couts[i,j] * x[i,j]
                for i in usines for j in entrepots)
\end{lstlisting}

\subsection{Contraintes}

Les contraintes de capacité des usines sont :
\[
x_{11} + x_{12} + x_{13} \leq 100
\]
\[
x_{21} + x_{22} + x_{23} \leq 150
\]

Les contraintes de demande des entrepôts sont :
\[
x_{11} + x_{21} = 80
\]
\[
x_{12} + x_{22} = 120
\]
\[
x_{13} + x_{23} = 50
\]

\subsection{Résultat obtenu}

\begin{lstlisting}
Usine 1 -> Entrepot 1 : 50 unites
Usine 1 -> Entrepot 3 : 50 unites
Usine 2 -> Entrepot 1 : 30 unites
Usine 2 -> Entrepot 2 : 120 unites

Cout total minimum = 780 MRU
\end{lstlisting}

La solution optimale permet de satisfaire toutes les demandes
avec un coût total minimal de \textbf{780 MRU}.

%==============================================
\section{Transport de Poissons -- Mauritanie}
%==============================================

\subsection{Description}
Le fichier \texttt{poisson.py} traite le transport de poissons
frais entre Nouadhibou et les villes de consommation.
L'objectif est d'analyser :
\begin{itemize}
  \item le coût minimal de transport ;
  \item le respect de la contrainte de périssabilité (12h max) ;
  \item la satisfaction des demandes de chaque ville.
\end{itemize}

\subsection{Modèle mathématique}

Le modèle utilisé est le modèle de transport avec contrainte
de temps. Le taux de périssabilité est donné par :
\[
x_{pv} = 0 \quad \text{si } t_{pv} > T_{\max} = 12 \text{ heures}
\]

avec :
\begin{itemize}
  \item $p$ : port d'origine ;
  \item $v$ : ville de destination ;
  \item $t_{pv}$ : temps de trajet en heures.
\end{itemize}

\subsection{Extrait du code}

\begin{lstlisting}
for p in ports:
    for v in villes:
        if temps[p,v] > temps_max:
            modele += x[p,v] == 0
\end{lstlisting}

Dans notre cas :
\[
T_{\max} = 12 \text{ heures}
\]

Comme tous les trajets sont inférieurs à 12h, le système
est faisable.

\subsection{Résultat obtenu}

\begin{lstlisting}
Tanit -> Rosso       : 80 tonnes (6h)  -> 30 400 MRU

Verification des demandes :
Nouakchott : 150/150 tonnes
Rosso      : 80/80 tonnes
Atar       : 100/100 tonnes
\end{lstlisting}

Ces résultats montrent que toutes les demandes sont satisfaites
dans les délais requis.

%==============================================
\section{Conclusion}
%==============================================

Ce projet nous a permis d'appliquer trois méthodes de Recherche
Opérationnelle avec Python.

Nous avons étudié :
\begin{itemize}
  \item la programmation en nombres entiers pour le Knapsack ;
  \item la programmation linéaire pour le transport ;
  \item un modèle réel mauritanien pour le transport de poissons.
\end{itemize}

L'utilisation de Python et PuLP facilite la modélisation, le
calcul et l'interprétation des résultats. Ce travail montre
l'importance de la Recherche Opérationnelle dans la logistique,
la santé et l'industrie.

\end{document}